% Shared implementation appendix. Both active papers include this file in full.
\section{Runtime and Initial Harness}
\label{app:surfaces}

The implementation described here is the repository at revision
\texttt{070d1a8493d7}, with the OOLONG-Pairs configuration inspected on
September 28, 2026. These are implementation specifications for the forthcoming
study, not experimental results. The study's eventual launch revision, frozen
configuration, split manifest, and initial/final harness hashes must accompany
its results.

\paragraph{Initialization.}
\hzero{} supplies a minimal REPL contract and generic instruction strings, a
disabled runtime policy, the default metadata formatter, empty helper and skill
libraries, and accepting answer middleware. \hzerostar{} replaces S1 with the
upstream-style system prompt, leaves S2--S5 empty, and enables the orchestrator
addendum. \texttt{H0*R}, the configured initial harness, additionally replaces
S1's recursive-call description and supplies a hand-authored S2 example. That
example distinguishes a bare completion from a child with its own REPL, splits
line-oriented input into batches of 300, requests label counts, merges counts in
Python, and retries a chunk after a parse failure. It was informed by development
observations; gains over \hzero{} alone would conflate this initialization with
automated editing. The \texttt{initial} evaluation condition loads the exact
saved round-one harness rather than substituting a reference from today's
registry. Figure~\ref{fig:edit-surfaces} and Figure~\ref{fig:full-repl-contract}
show the sparse reference, not the complete \texttt{H0*R} prompt.
At this revision the registry \texttt{H0*R} serialization hash is
\path{0ca3510ded8a0ce5095c2e3338c4d8606fbc819650c20b65079b405ae1ecac2b}.
The study's saved initial artifact, including its rendered prompt and runtime
settings, is the authoritative comparator.

\paragraph{Execution contract.}
S1--S5 are concatenated into one system prompt; a clause referring to early turns
is guidance, not a separately timed host hook. Recursive \texttt{rlm\_query}
children below the depth limit receive their own REPL and the shared prompt.
Bare \texttt{llm\_query} calls and leaves at the recursion limit are completions,
without that REPL. Printed code output joins accumulated message history;
input-as-variable does not prevent code from printing input records. The normal
answer path uses the REPL answer dictionary, while iteration exhaustion can
invoke the runtime's final-answer fallback. Admission checks preserve specified
interfaces and caps, not a formal proof of arbitrary generated code's behavior.

\begin{table}[htbp]
\centering\small
\renewcommand{\arraystretch}{1.12}
\begin{tabularx}{\linewidth}{@{}p{0.19\linewidth}X@{}}
\toprule
Surface & Editable payload, visible information, and execution scope \\
\midrule
S1: REPL contract & Full text replacement describing the environment/API; shared root/recursive-child prompt. \\
S2: Decomposition & Full text replacement guiding input inspection and decomposition; shared prompt, scoped to the current subtask. \\
S3: Execution & Full text replacement guiding parsing, calls, aggregation, and task predicates; shared prompt. \\
S4: Verification & Full text replacement guiding checks before submission; shared prompt, with no independent semantic oracle. \\
S5: Recovery & Full text replacement guiding actions after visible errors; cannot extend a terminated run or add host retries. \\
S6: Runtime policy & Dictionary with enabled flag, prompt/batch caps, depth, syntax-retry/count, and output-validation settings. Enforcement is root local except inherited depth. Width/length caps refuse work; syntax retry repeats the same prompt. \\
S7: Metadata & Root formatter function receiving bounded stdout and redacted variable types/lengths; preserves its declared output bound. It cannot inspect hidden variable values or the full call history. \\
S8: REPL helpers & One named function added/replaced in one root or child helper dictionary. Helpers execute only when called; the proposer must identify a caller and input/output contract. \\
S9: Answer middleware & Root function receiving answer text and redacted inventory; may accept transformed text or redirect with a nudge. It cannot inspect hidden records or classifications. \\
S10: Skills & One named text procedure added, replaced, or removed. A shared name/description index advertises it; its body matters only when loaded or forwarded. \\
\bottomrule
\end{tabularx}
\caption{Ten surfaces and their actual capabilities. Provider/model selection,
decoding/reasoning settings, verifiers, and experiment-owned caps are outside the
edit space. S8 spans two fields; a surface is not necessarily one builder.}
\label{tab:surfaces}
\end{table}

\paragraph{Skills and token accounting.}
The S10 index enters the system prompt after S5. For a nonempty library the
runner installs \texttt{load\_skill(name)} in root and recursive-child REPLs.
The host merges a single-entry edit with the incumbent library, including
explicit removal. Limits are eight entries, 40 characters per name, 200 per
description, 4,000 per body, and 16,000 total. A loaded body stored only in a
variable need not be model visible. If printed or forwarded it incurs input
tokens, potentially on later turns while it remains in history. Neither loading
a body nor advertising it proves that its procedure was followed.

% Figure 1: the initial harness H0 and its ten editable surfaces.
% Body only; the `harness' listing style, its colours, and \hzero live in the
% preamble of proposal.tex, matching how Figure_2.tex relies on the tikz styling
% defined there.
\begin{figure}[!tbp]
    \centering
    \begin{lstlisting}[style=harness]
    # shrlm/rlm_harness.py: the ten surfaces Self-Harness may edit (S1 to S10).
    
    # Condensed for visibility; the full literal appears in the appendix.
    MINIMAL_REPL_CONTRACT = """
    You are an RLM operating through a persistent Python REPL.
    Use `context`, flat or recursive LM calls, `SHOW_VARS()`, and `answer`.
    Execute one ```repl``` block per turn; print inspections and mark the
    answer ready only when complete. [Display-only condensation.]
    """
    
    def build_repl_contract() -> str:                                    # S1
        return MINIMAL_REPL_CONTRACT
    
    def build_decomposition_instruction() -> str:                        # S2
        return ("Start by probing `context` to understand what you have, "
                "then decide how the task breaks into steps.")
    
    def build_execution_instruction() -> str:                            # S3
        return ("On each turn, run one block of code, print what you need "
                "to see, and build the result up in REPL variables.")
    
    def build_verification_instruction() -> str:                         # S4
        return 'Check your candidate answer before setting `answer["ready"] = True`.'
    
    def build_recovery_instruction() -> str:                             # S5
        return ("If a sub-call errors or returns something you cannot use, "
                "adjust the approach and try again.")
    
    def build_runtime_policy() -> dict[str, Any]:                        # S6
        return {
            "enabled": False,
            "max_prompt_chars": None,
            "max_batch_width": None,
            "max_depth": None,
            "retry_on_syntax_error": None,
            "max_retries": None,
            "validate_sub_output": None,
        }
    
    def build_metadata(                                                  # S7
        stdout: str,
        repl_inventory: dict[str, tuple[str, int]],
        max_character_length: int = DEFAULT_MAX_CHARACTER_LENGTH,
    ) -> str:
        return default_metadata_builder(stdout, repl_inventory,
                                        max_character_length)
    build_metadata.declared_bound = DEFAULT_MAX_CHARACTER_LENGTH        # S7
    
    def build_repl_helpers() -> dict[str, Any]:                          # S8
        return {}
    
    def build_sub_repl_helpers() -> dict[str, Any]:                      # S8
        return {}
    
    def build_answer_middleware() -> AnswerMiddlewareFn:                 # S9
        return accept_answer
    
    @dataclass(frozen=True)
    class SkillEntry:                                                    # S10
        name: str          # REPL-safe identifier; index label
        description: str   # one line: when to consult it; enters the prompt
        body: str          # ordered steps; served verbatim by `load_skill(name)`
    
    def build_skills() -> list[SkillEntry]:                              # S10
        return []
    \end{lstlisting}
    \caption{Initial harness \hzero{}, the starting point for \sh{}: ten declared
    surfaces, one builder function each, keyed to phases of the \rlm{} turn loop.
  %  Eight of the ten return nothing, a disabled policy, or a single generic line, so the orchestration strategy a tuned harness carries must be recovered by the optimization loop from its own execution traces. 
    The \sid{S1} contract is
    condensed here only to keep the complete interface visible on one page;
    \cref{fig:full-repl-contract} reproduces the exact initial prompt. \sid{S10} is the skill library:
    named procedures, each a \texttt{name}, a one-line \texttt{description} of when to
    consult it, and a \texttt{body} of ordered steps. Only the name/description index
    enters the system prompt, after \sid{S5}; a body is returned verbatim by
    \texttt{load\_skill(name)}, which the runner installs in the root and child REPLs
    when the list is non-empty, so a procedure costs tokens only on the turn that loads
    it. At most eight entries; \hzero{} carries none. \sid{S6}, \sid{S7} and \sid{S9}
    required a runtime injection point in the reference implementation of
    \citet{zhang2025rlm}; \sid{S10} adds one hook to record skill loads in the trace;
    the rest use its existing prompt and tool configuration. A proposal modifies
    exactly one surface.}
    \label{fig:edit-surfaces}
\end{figure}
    

% Appendix figure: exact, unabridged sparse H0 S1 template (not H0*R).
% Figure 1 intentionally contains a display-only condensation of this prompt.
\begin{figure}[!tbp]
    \centering
    \begin{lstlisting}[style=harness]
    MINIMAL_REPL_CONTRACT = textwrap.dedent(
        """\
        You are a Recursive Language Model (RLM): a language model whose prompt-related information lives in a Python REPL that you drive turn by turn. You will be queried turn-by-turn until you have an answer to the query.

        Available in the REPL:
        - `context`: the information related to the prompt (typically `str` or `list[str]`).
        - `llm_query(prompt: str, model: str | None = None) -> str`: a single sub-LLM completion.
        - `llm_query_batched(prompts: list[str], model=None) -> list[str]`: several sub-LLM completions in parallel; same order out as in.
        - `rlm_query(prompt, model=None)` / `rlm_query_batched(prompts, model=None)`: recursive RLM sub-calls. Fall back to `llm_query` / `llm_query_batched` when recursion is disabled.
        - `SHOW_VARS() -> str`: list every variable currently in the REPL.
        - `answer`: dict initialized to `{{"content": "", "ready": False}}`. To submit, set `answer["content"]` to the final answer and `answer["ready"] = True` inside a ```repl``` block.
        {custom_tools_section}

        Write code in ```repl``` blocks, one block per turn; the REPL persists across turns. The REPL is NOT a Jupyter cell - only `print(...)` output (stdout) is shown back to you between turns; a bare expression on the last line is silently discarded. REPL outputs over ~20K characters are truncated."""
    )
    \end{lstlisting}
    \caption{The complete sparse \hzero{} S1 source template, before placeholder
    formatting. Figure~\ref{fig:edit-surfaces} uses a display-only condensation.
    This is a reference contract, not the full initial prompt of the configured
    \texttt{H0*R} study; its upstream-style prompt, recursive-call change, S2
    example, and orchestrator addendum are specified in
    Appendix~\ref{app:surfaces}.}
    \label{fig:full-repl-contract}
\end{figure}


\section{Mining, Diagnosis, and Evidence Selection}
\label{app:mining}

Final-answer verification supplies pass/fail and an outcome category. Eligible
failed runs are sent to a fixed-model attributor; passing runs and configured
environment-owned exclusions are not sent through failure attribution. Excluded,
unattributed, and rejected records remain visible in accounting. Each diagnosis
records an inferred mechanism, location and causal status, supporting operation
references, and uncertainty. Structural absence of recursion can be observed
directly; semantic correctness and root-versus-child causality are not implied
by final-answer scoring.

The attributor must distinguish skipped records, wrong or uncertain labels, and
incorrect aggregation or predicates. Missing answer elements alone do not prove
input loss. It must account for observed checks and repairs, and distinguish
missing evidence from evidence that an action was omitted. For
\texttt{incomplete\_coverage}, a structured basis records status, input scope,
loss observation, and counterevidence. Statuses \texttt{not\_established} and
\texttt{contradicted} normalize to \texttt{OTHER}/\texttt{UNATTRIBUTED}, rather
than authorizing a coverage edit. Host acceptance verifies schema, vocabulary,
field bounds, and resolvable references, not the truth of a causal explanation.
Attribution response attempts default to three; transport retries are separate.

\paragraph{Clustering and routing.}
Records cluster by exact agreement on
\[
\phi(r)=(\text{outcome cause},\ \text{recorded failing level},\
         \text{causal status},\ \text{mechanism}).
\]
The signature mixes observed outcomes with model-inferred fields; it is not a
verified causal label. Patterns rank first by distinct-instance support, then
actionability, then a deterministic signature tie-break. Attempt support is
recorded separately. Minimum support two is a flag, not automatic deletion.
Three representatives and a fixed actionability heuristic guide reading order:
weights are 0.4 causal status, 0.3 grounding availability, 0.2 known mechanism,
and 0.1 normalized-description homogeneity. The last term uses distinct text
strings, not semantic embeddings; the score is not a calibrated success
probability.

\begin{table}[htbp]
\centering\small
\begin{tabularx}{\linewidth}{@{}Xl@{}}
\toprule
Mechanism & Eligible surfaces (primary first) \\
\midrule
REPL contract misuse & S1, S9, S10, S4 \\
Whole-input subcall collapse & S2, S3, S10, S6 \\
Incomplete coverage & S2, S4, S3, S10 \\
Redundant decomposition & S2, S3, S10 \\
Misaligned unit & S2, S10 \\
Depth degradation & S3, S6, S2 \\
Insufficient recursion & S3, S2, S6 \\
Bare completion substituted for recursion & S3, S1, S10 \\
Premature termination & S4, S9, S3 \\
Skipped verification & S4, S9, S10 \\
Swallowed subcall error & S5, S3, S6 \\
Iteration budget exhaustion & S6, S3, S1, S10, S5$^{\dagger}$ \\
Unparsed child output & S7, S3, S5, S8$^{\dagger}$ \\
REPL execution fault & S8, S10, S3, S1, S5$^{\dagger}$ \\
Lossy aggregation & S3, S4, S8$^{\dagger}$ \\
Unconsulted procedure & S10, S3, S2 \\
Other & S1--S10 \\
\bottomrule
\end{tabularx}
\caption{Human-defined, many-to-many routing (taxonomy 3.3.0).
An \texttt{UNATTRIBUTED} diagnosis has no eligible route, including when its
mechanism is Other. $\dagger$ denotes additional operation-evidence requirements.}
\label{tab:routes}
\end{table}

Conditional S8 eligibility requires a complete cited operation; unparsed child
output additionally requires child-response evidence. Conditional S5 eligibility
requires an observed error and a subsequent complete operation. These are
structural opportunities to intervene, not proof of successful recovery or an
effective proposed fix. Aggregation errors route to S3/S4 or supported S8;
S9 remains limited to answer-visible properties.

\paragraph{Bounded evidence.}
The attribution digest defaults to 12,000 characters. Proposal evidence defaults
to 32,000 characters, separately from the 12,000-character history budget.
These are component budgets, not a total prompt cap. A compact pattern inventory
is supplemented by up to $\min(k,4)$ expanded patterns, preferring distinct
actionable mechanisms before additional signatures/context. Inventory-only
patterns do not authorize proposals. Evidence follows cited child calls and
relevant parsing, merging, or filtering operations, with nearby input definitions,
consumers, and corrections. Whole code operations are retained where they fit;
required operations that cannot fit are omitted with scope markers rather than
presented as complete fragments. Outputs and prose may be shortened. A passing
contrast sharing operation names can be included when available. This improves
context for comparison but guarantees neither full causal context nor balanced
surface coverage.

\section{Proposal, Materialization, and History}
\label{app:proposal}

One response under \texttt{proposal-selection/v2} selects evidence-supported
pattern/surface interventions before supplying replacement payloads; these are
not separate selection and generation calls. The first admitted candidate owns
its surface and pattern. Rejected entries do not reserve a surface. At most
$k$ edits are admitted, with fewer or zero allowed. Three short explanation
fields identify incumbent behavior, unresolved observed failure, and the change
at a specific operation, input, or trigger. Guidance asks whether the same
demonstrated failure could persist if the edit were followed perfectly. Repeating
an existing instruction more emphatically, moving it to another surface, or
repairing an already-recovered fault is insufficient without a substantive
behavioral difference. Replacing a misleading example can be the smallest
effective change.

The general prompt asks which information must survive each operation and which
task conditions must the final computation enforce. It does not prescribe an
OOLONG-specific record-ID algorithm. Environment-provided API and answer
contracts remain authoritative. These instructions guide model judgment; they
are not deterministic tests of semantic novelty or correctness.

\paragraph{Repair and preflight.}
Response/schema attempts are bounded (default maximum eight), distinct from
task validations. When failed-member repair is invoked there is one repair
response, restricted to failed original patterns. It may select another eligible
unoccupied surface or withdraw an edit while retaining valid siblings unchanged;
it may not invent new patterns. Collision repair must choose a contender or
withdraw it, still enforcing one admitted edit per surface.

For text under \texttt{literal-text/v1}, the model supplies literal instructions;
the host encodes formatting braces exactly once and restores the designated
custom-tools placeholder. S1--S5 replace full text, S6 replaces a policy
dictionary, S7/S9 supply functions with required interfaces, S8 edits one named
helper in one dictionary, and S10 edits one named library entry. Preflight checks
serialization and base hashes, an actual single-surface change, reserved
interfaces, caps, formatting, and selected smoke/answer fixtures. Environment
answer fixtures are synthetic contract examples, not validation tasks. Literal
no-op rejection does not detect every behaviorally redundant instruction.
Saved artifacts distinguish literal proposal text, encoded source/template, and
runtime rendering.

\paragraph{History and qualified progress.}
The archive retains attempts, admission/rejection reasons, batch membership,
measured outcomes, and revision links. The bounded prompt summary prioritizes
recent rounds, the incumbent's promotion, relevant measured predecessors, and an
older promising direction when it fits. Revisiting an intervention requires a
resolvable predecessor reference and a stated substantive change to the operation
or joint hypothesis. Validation contributes aggregate scores, costs, sanitized
failure summaries, and observed activation; its raw traces and gold/answer dumps
are not forwarded to the proposer.

A rejected measured candidate can be marked potentially promising when a
comparable verifier-defined dense metric improves in its declared direction.
Definitions specify identity/version, direction, aggregation, precision, and
terminal-outcome values. Comparison also requires compatible verifier/protocol/
repetition settings and matching instance-attempt sets. Missing, unsupported, or
incomparable metrics remain \texttt{not\_assessed}. The label is qualified
evidence, not promotion or statistical significance; a simultaneous exact-pass
regression remains recorded. An identical rejected batch already measured under
the same incumbent can be skipped before new validation.

\paragraph{Observed behavior.}
Instrumentation reports root S6 syntax retries and S9 redirects, and S10 loads
across the traced call tree, with observation coverage. Missing coverage is
unassessed. The intended-behavior field remains \texttt{not\_assessed}: activation
does not establish compliance or task correctness. There is no general semantic
compliance detector for S1--S5, S7, or S8.

\section{Joint Validation and Stopping}
\label{app:validation}
\label{app:algorithm}

Let $H$ be the incumbent, $C$ the composition of all admitted edits, and
$P(X)$ the number of final-answer verifier passes from harness $X$ on the same
$n_{\mathrm{ho}}v$ validation attempts. A singleton is the combined candidate;
an empty or skipped batch schedules no validation. Only $H$ and $C$ are scored;
constituents of a multi-edit batch have membership records without individual
scores. The implemented pass-count conditions are
\begin{equation}
\Delta=P(C)-P(H),\qquad
\Delta\geq-\tau_{\mathrm{reg}},\qquad
\Delta\geq\tau_{\mathrm{imp}}.
\label{eq:promotion}
\end{equation}
For each configured resource metric $q$, its mean must lie within the inclusive
baseline-relative band, $a_q\bar q(H)\leq\bar q(C)\leq b_q\bar q(H)$, with
unbounded endpoints permitted. Current thresholds are both zero; mean cost must
be between zero and three times incumbent cost, and subcall count is unconstrained.
Thus an exact-pass tie, including zero versus zero, can promote even if F1 falls.
This is the intended aggregate verifier-pass/resource rule, not a guarantee of
per-task improvement or secondary-metric non-regression.

\begin{quote}\small
\textbf{Inputs:} configured initial harness $H$, mining tasks $M$, validation
tasks $V$, repetitions $m,v$, edit ceiling $k$, round ceiling $T$, patience $p$,
and promotion settings. Initialize history and no-promotion counter $z=0$.
Repeat steps 1--4 for at most $T$ rounds; step 5 follows loop termination.
\begin{enumerate}
\item Execute or resume $H$ on $M$ with $m$ attempts;
record final outcomes, traces, and available observations.
\item Diagnose eligible failures, cluster records, and select bounded evidence
and history. Propose, admit, and if needed repair up to $k$ edits on distinct
patterns and surfaces.
\item Set promoted to false. If an evaluable, nonduplicate batch exists, compose
$C$ and freshly evaluate or resume $H,C$ on $V$ with $v$ attempts. An unscorable
baseline follows the experiment error path. For a complete comparable candidate,
apply Equation~\ref{eq:promotion} and resource bands; if accepted set $H=C$ and
promoted to true.
\item Persist the batch outcome and constituent membership, or the reason no
comparison occurred; update history. Set $z=0$ on promotion, otherwise $z=z+1$.
Stop when $z\geq p$.
\item Freeze the last incumbent; run separately specified final-evaluation
conditions. Do not retrospectively select a harness using test scores or F1.
\end{enumerate}
\end{quote}

\paragraph{Failures, denominators, and resume.}
Candidate-owned execution errors can become failed attempts without terminating
the whole experiment. Configuration/operational/provider errors, cancellation,
memory or I/O failures, and observation-persistence failures may still propagate.
Complete comparisons require matching planned attempts; a partial baseline cannot
anchor promotion. The final-evaluation driver can retain completed sibling
results while reporting an incomplete or over-budget condition.

\begin{table}[htbp]
\centering\small
\begin{tabularx}{\linewidth}{@{}p{0.27\linewidth}X@{}}
\toprule
Recorded status & Treatment \\
\midrule
Completed correct/incorrect & Enters attempt denominator; only a verifier pass counts as success. \\
Completed execution/resource/format/filter failure & Enters denominator as a failure; dense terminal values follow the verifier definition. Partial usage/trace coverage remains labeled. \\
Missing/canceled attempt & Report separately, not as a fabricated measured zero. Incomplete validation cannot silently become a complete comparison. \\
Unscorable baseline & Persist evidence and stop/raise through the experiment error path; this is not a candidate rejection. \\
Resumed/cached completion & Reuse only with matching identity/contracts; not a new independent attempt or new spend. \\
\bottomrule
\end{tabularx}
\caption{Outcome accounting for optimization and final evaluation.}
\end{table}

\section{Experimental Protocol and Provenance}
\label{app:protocol}

\paragraph{Pinned data and scoring.}
The profile is \path{configs/experiment_oolong_pairs_DeepSeekV4Flash.toml}.
The dataset is \texttt{oolongbench/oolong-synth}, upstream validation partition,
\texttt{trec\_coarse}, revision
\path{f0d59eaf0febf130664cfceb710436c8e3216b2b}. The loader selects the first
two distinct windows per configured length, with a 50,000-row scan ceiling,
then pairs each with task IDs 1--20 and shuffles the instances with seed zero.
The seed controls task partitioning, not per-attempt model randomness. Labels
are removed from model-visible records; gold user-pair sets are computed from
the labeled records and task predicates.

The scorer extracts integer pairs and compares sets. Global pair order and
duplicate pairs do not affect set equality, but order within a pair is preserved:
the lower user ID must come first. A bracketed list or newline-separated pairs
are accepted. The explicit empty marker is ``No valid pairs found.''; literal
\texttt{[]} is not an empty answer. No parseable pair or marker is a format
failure. A correctly empty prediction against an empty gold set passes and has
precision/recall/F1 equal to one. These extraction and empty-set distinctions
are adaptations of the example scorer, not an unmodified upstream verifier.

\paragraph{Task separation and shared records.}
The two short windows contain 188 records each; the two long windows contain
6,374 each. Short instances are partitioned 20/10/10, while all 40 long
instances are reserved for separate evaluation. The current split remains
unchanged. Table~\ref{tab:context-overlap} makes its unit explicit.

\begin{table}[htbp]
\centering\small
\begin{tabular}{@{}lrrr@{}}
\toprule
Short context (input-prefix hash) & Mining & Validation & Short test \\
\midrule
Window 9 (\texttt{5f067945b676}) & 10 & 6 & 4 \\
Window 10 (\texttt{e939371361d4}) & 10 & 4 & 6 \\
\bottomrule
\end{tabular}
\caption{Task counts per shared short record collection at partition seed zero.
Hashes exclude the task question. Distinct task instances do not imply distinct
underlying record collections.}
\label{tab:context-overlap}
\end{table}

Different predicates over the same window ask different aggregation questions
over identical records. Between the two short windows, no complete
\texttt{(user\_id, date, question\_text)} record is identical, but 43 question
texts recur with different users/dates. The long windows also have no identical
complete-record overlap with the short windows, while sharing question texts.
Short-test results support the held-out-task-instance interpretation; long
evaluation assesses larger record collections. Neither establishes wholly
unseen-source-text generalization. A stronger source-separated design is needed
only if that becomes a central claim. This overlap is not evidence of answer
memorization.

The nominal upstream lengths differ by $32\times$, but stripping labels and
adding task text changes actual prompt lengths. Evaluation will report
reconstructed prompt character counts and, when measured, token counts with
the tokenizer identified, separately from nominal window labels. Multiple
predicates per window and three requests per task do not increase the number of
independent contexts. Any future interval or paired test must specify its unit,
attempt aggregation, missingness, and the two-context limitation; no implemented
bootstrap or significance analysis is claimed here.

\paragraph{Model, caps, and concurrency.}
All three optimization roles use Azure Foundry \texttt{DeepSeek-V4-Flash},
temperature 0, top-$p$ 0.95, and an output-token cap of 8,192. The provider's
default reasoning mode is enabled, without guaranteeing returned reasoning.
Per-run limits are 30 iterations, depth three, 3,600 seconds, and a \$2 usage
budget. Each validation subject, including the incumbent, has a \$100 cumulative
budget. There are three mining run workers. Validation allows five subject
workers but has only two measured subjects, with one run worker per subject;
recursive calls can add provider concurrency. Final evaluation has a separate
run-worker option and uses three attempts per task per condition. Record the
chosen worker count in the evaluation manifest; it does not change repetition
counts.

The main comparison explicitly selects \texttt{initial} and \texttt{sh\_rlm}.
The default reference grid is \texttt{b1}, \texttt{h0\_star},
\texttt{lambda\_rlm}, \texttt{sh\_rlm}. The latter's typed baseline is identified
by its method kind, reconstruction version, upstream revision, and source hash;
it is not an unmodified released OOLONG-Pairs implementation. Model/deployment,
rendered prompts, source/configuration snapshots, split IDs/hashes, exact edit
payloads, harness lineage, and analysis versions belong in the study artifacts.
Freezing these before a run does not by itself constitute preregistration.
Previously inspected development/test material must retain its exposure history
if later used to support a result.

\section{Measurement, Cost, and Saved Observations}
\label{app:measurement}
\label{sec:cost}

\paragraph{Quality and surface activity.}
Exact pass is set equality, not rounded F1 equal to one. OOLONG-Pairs stores
per-attempt precision, recall, and F1 to three decimals; reported mean F1 averages
recorded attempt values, including verifier-declared zeros for completed runtime,
resource, format, and content-filter failures. It is macro averaging across
attempts, not pooled pair-level micro-F1. Missing/extra counts are reported only
where available, with coverage, rather than filling absent diagnostics with zero.
Other environments supply their own metric definitions or remain unassessed.

Surface activity distinguishes submitted, admitted, evaluated, and promoted
edits. Constituents inherit promotion membership from a promoted batch without
individual efficacy estimates. Per-round quality plots use the freshly measured
candidate on promotion and freshly measured incumbent otherwise; a round without
a measurement is missing, not a zero or a carried-forward score. Environment
labels come from the configuration. Trace summaries can report depth, call
counts, and the implemented whole-input-subcall proxy: maximum descendant prompt
characters divided by root prompt characters exceeds 0.8. Its denominator is
walkable runs, passing and failing, with unwalkable runs separate; it is neither
a token ratio nor proof of meaningless decomposition.

\paragraph{Cost.}
An ordinary completed round with an evaluable batch schedules
\begin{equation}
N_{\mathrm{round}}=m n_{\mathrm{in}}+2v n_{\mathrm{ho}}.
\label{eq:round-cost}
\end{equation}
The current profile gives $40+20=60$ task attempts, independent of how many
edits the batch contains, or 180 over three such rounds. No evaluable batch
means no new validation leg. With 10 short and 40 long tasks, final evaluation
requires $3(10+40)J=150J$ attempts for $J$ conditions: 300 for initial versus
final, 600 for the default four-condition grid, or 750 for all five. These are
scheduling counts, not measured results or provider-call counts. Root attempts
may issue many recursive/model requests. Attribution, proposal, repairs, and
operational retries also consume resources.

Configured accounting rates are \$0.19 per million input tokens and \$0.51 per
million output tokens, as recorded for this study; they are not a current market
quote. Actual usage and its coverage will be reported separately from scheduling
projections. Task usage is not added again to a stage total
that already includes it. Breaker reservations are not billed spend; interrupted
or unreported requests can leave a lower bound, and the per-run budget is not a
guarantee of complete invoice accounting.

\paragraph{Provider-returned observations.}
Instrumented execution, mining, validation, final evaluation, attribution,
proposal/repair, recursive calls, compaction, and fallback paths save returned
response/reasoning observations separately from executable code and final
answers. References carry call IDs, paths, and hashes. Run traces point to
\path{<run_id>.llm_calls/<call_id>/}; attribution and proposal artifacts point to
their stage-specific \path{llm_calls/} directories. Started, response, and finished
events distinguish failed or interrupted calls. Cache hits retain source
provenance and do not add spend.

Availability is explicit: \texttt{returned}, \texttt{opaque\_only},
\texttt{not\_returned}, \texttt{unsupported\_client},
\texttt{legacy\_unavailable}, or \texttt{no\_response}. Returned content can be
text or a summary; this is not guaranteed access to full private reasoning.
Unknown reasoning-token counts stay unknown and are not added to reported output
usage again. Legacy traces cannot be backfilled with unreturned reasoning.
Observation artifacts are not automatically fed into attribution digests,
proposal evidence, or history. Ordinary assistant prose in a digest is a
different source. Observation persistence errors remain visible failures.

\begin{table}[htbp]
\centering\small
\begin{tabularx}{\linewidth}{@{}Xl@{}}
\toprule
Contract & Version \\
\midrule
Taxonomy / capabilities & 3.3.0 / \texttt{surface-capabilities/v1} \\
Attribution prompt / validator / digest & 1.6.0 / 1.2.0 / 1.8.0 \\
Proposal prompt / validator & 4.3.0 / 4.2.0 \\
Proposal response / text & \texttt{proposal-selection/v2} / \texttt{literal-text/v1} \\
Evidence selector / diagnostic history & 4.3.0 / 2.1.0 \\
Compact history & \texttt{proposal-history/v3} \\
Validation / summary & \texttt{heldout-batch/v2} / \texttt{shrlm-validation-summary/v3} \\
Behavior / provider observations & \texttt{behavior-observations/v1} / \texttt{rlm-llm-observation/v1} \\
\bottomrule
\end{tabularx}
\caption{Contracts at the documented revision; archive the versions actually
used by the forthcoming run with its results.}
\end{table}

\section{Optional Follow-up Studies}
\label{app:extensions}
\label{app:finetune}
\label{app:alignment}

Cross-environment transfer would require a separately specified target protocol
and completed evaluation. A weight-training comparison would require an available
trainable checkpoint, justified backbone matching, and measured compute; no
fine-tuning condition or hardware allocation is claimed for this hosted study.
Joint batch outcomes suffice for the present design. An optional future
leave-one-surface-out study would revert a final surface to its starting value
and rerun the same evaluation; it would measure that final-state counterfactual,
not the marginal contribution of every historical edit overwritten along the way.
No such ablation has been completed here.

An optional safety probe could compare frozen harness states for changes in
refusal and harmful-answer behavior, motivated by work on
misevolution~\citep{shao2025misevolution}. The task gate does not measure these
outcomes, so improved benchmark behavior would not establish safe deployment.
Any probe would need its own data, scoring, and budget. Textual reasoning
instructions are editable; provider reasoning settings are experiment owned.
